\subsection*{调试 RTTI 与 \texttt{debugger}}

当 \co@print@ 和 \co@println@ 不足以定位问题时，可以使用 Croqtile 提供的调试运行时类型信息，让 \co@debugger@ 能够查看 Croqtile 类型。

编译时使用 \co@-g@ 启用调试符号：

\begin{lstlisting}[language=croqtile,caption={生成带调试符号的程序}]
croq kernel.co -g -o kernel_debug
\end{lstlisting}

生成代码中包含用于描述 Croqtile 类型的 RTTI 结构体：

\begin{table}[htbp]\centering\small
\begin{tabularx}{\linewidth}{@{}lXX@{}}
\toprule
\textbf{Croqtile 类型} & \textbf{调试器中的类型} & \textbf{字段} \\
\midrule
带形状张量（\texttt{s32 [M, N]}） & \texttt{croq::rtti::spanned\textless{}int, 2\textgreater{}} & \texttt{.span.data[]}：各维大小；\texttt{.stride.data[]}：步长；\texttt{.data}：数据指针 \\
带范围的索引元组 & \texttt{croq::rtti::bounded\_ituple\textless{}N\textgreater{}} & \texttt{.data[]}：索引值；\texttt{.ub[]}：上界 \\
整数元组 & \texttt{croq::rtti::ituple\textless{}N\textgreater{}} & \texttt{.data[]}：整数值 \\
多维形状 & \texttt{croq::rtti::mdspan\textless{}N\textgreater{}} & \texttt{.data[]}：各维大小 \\
\bottomrule
\end{tabularx}
\end{table}

\begin{lstlisting}[language=croqtile,caption={在调试器中查看变量}]
debugger -q ./kernel_debug
(gdb) break my_kernel
(gdb) run
(gdb) ptype __dbg_lhs
type = struct croq::rtti::spanned<int, 2>
(gdb) print __dbg_lhs.span.data[0]
$1 = 32
(gdb) print __dbg_lhs.span.data[1]
$2 = 64
(gdb) print __dbg_lhs.data != 0
$3 = true
\end{lstlisting}

变量名中的 \co@__dbg_@ 前缀由编译器生成，使 Croqtile 变量能够与生成的 C++ 中间变量一起出现在调试器中。
